% Copyright Xucong Lyu
% Version 1.0
% 本模板基于《河南师范大学研究生学位论文格式要求（2026年更新）》制作
% 本模板适用于硕士及博士研究生毕业论文的写作
% 编译：xelatex

% 选项 science（理工）或 humanities（人文社科）
\documentclass[science]{hnnuthesis}

% 论文版本开关，注释掉%\blindreviewtrue ，论文切换为完整版，不注释为盲审版
% 盲审版隐去了作者姓名、导师姓名、学号，以及致谢、发表文章目录和独创性声明等
\newif\ifblindreview
%\blindreviewtrue         % 盲审版


%以下封面及论文信息部分
\providecommand{\HNNUtitlebreak}{}
\hnnusetup{title={题目题目题目题目题目题目题目
                  \HNNUtitlebreak%题目过长使用此代码换行
                   题目题目题目},
           englishtitle={title title title title title title title title title title title title},
           author={张\quad 三},
           studentid={2XXXXXXXX},
           classification={Q178},
           major={某某、某某},
           direction={某某方向},
           degree={（学术学位）},
           degreecategory={农学博士},
           level={博士},
           supervisor={李\quad 四\quad 教授},
           date={二〇二七年五月},
           dateArabic={2027年5月},
           keywords={学位论文；排版规范；LaTeX},
           enkeywords={thesis, typesetting, LaTeX},
englishdegree={Doctor of Agriculture},%英文封面的学位说明
englishauthor={Zhang San},
englishsupervisor={Prof. Li Si},
englishdate={June, 2027}}

% 编译盲审需要隐去的部分
\ifblindreview
  \renewcommand{\HNNUauthor}{}
  \renewcommand{\HNNUsupervisor}{}
  \renewcommand{\HNNUenglishauthor}{}
  \renewcommand{\HNNUenglishsupervisor}{}
  \renewcommand{\HNNUstudentid}{}
  \hypersetup{pdfauthor={}}
\fi


% 以下是论文主体部分
\begin{document}
\hnnumakecover
\hnnutitlepage
\hnnufrontmatter

% 中文摘要
\begin{hnnuabstract}
本文以河南师范大学研究生学位论文格式要求为依据，构建可重复使用的学位论文排版模板。研究方法包括规范分析、样式设计和编译验证。此处为演示文本，正式提交时应替换为完整摘要。
\end{hnnuabstract}

% 英文摘要
\begin{hnnuenglishabstract}
This sample demonstrates a thesis template based on the graduate thesis formatting requirements of Henan Normal University. Replace this paragraph with your actual English abstract corresponding to the Chinese version.
\end{hnnuenglishabstract}

\cleardoublepage
\phantomsection
\addcontentsline{toc}{chapter}{目\quad 录}
\hnnutableofcontents
% 图表较多时启用以下清单：
%\listoffigures
%\cleardoublepage
%\listoftables
%\cleardoublepage
%\hnnuchapter{主要符号表}
%\noindent $x$：自变量；$y$：因变量。
\hnnumainmatter

% 以下为论文正文
% 由于毕业论文篇幅较长，各小节分开编写较好，最后引入主体
\chapter{绪论}
\section{研究背景}
此处为研究背景。此处为研究背景。此处为研究背景。此处为研究背景。此处为研究背景。此处为研究背景。此处为研究背景\cite{knuth1984,lamport1994}。
\subsection{研究意义}
论文正文采用小四号宋体\cite{lyu_environmental_2026}、首行缩进两字符及一点五倍行距\cite{lyu_environmental_2026}。
\section{研究内容}
图\ref{demo}为图题编号示例。

\begin{figure}[htbp]
	\centering
	\includegraphics[width=.5\textwidth]{hnnu-calligraphy.png}
	  \bicaption
        {河南师范大学}
        {Henan Normal University}
	\label{demo}
\end{figure}

\chapter{这里是第二章}
\section{引言}
这里是引言这里是引言

\section{材料与方法}
\subsection{研究区域}


\chapter{这里是第三章}
\section{引言}


\include{chapters/04}
\include{chapters/05}


\cleardoublepage
\phantomsection
\addcontentsline{toc}{chapter}{参考文献}
\bibliography{references}

% ========================================
% 以下内容仅最终版显示
% ========================================
\ifblindreview
  % 盲审版：参考文献后直接结束
\else
  % 如有附录，放在这里
  \appendix
  % 致谢
  \hnnuacknowledgements{%
    感谢导师在论文选题、研究与写作过程中给予的指导。
  }

  % 攻读学位期间发表的学术论文目录
  \hnnupublications{%
    \noindent [1] 示例论文信息。
  }

  % 独创性声明和论文使用授权说明
  \hnnudeclarations

\fi

\end{document}